\chapter{Primeros pasos}
\label{sec:man-quickstart}

Esta guía te lleva de la app recién instalada a un hunt con
resultados y scoring, sobre datos sintéticos, en unos diez
minutos.

\section{Datos demo}

El repositorio incluye \filepath{demo_data/}. Para esta guía usamos
\filepath{demo_data/fortigate_sample.csv}: ~50k líneas de
FortiAnalyzer sintético con un incidente de lateral movement
embebido.

Si querés datos reales, ver apéndice~\ref{app:man-demo}.

\section{1. Crear sesión}

Abrí la app. \emph{Archivo \textrightarrow{} Nueva sesión}. Dale
nombre (ej.\ \code{quickstart-forti}). Aparece el panel de ingesta
a la izquierda.

\section{2. Ingestar}

Arrastrá \filepath{demo_data/fortigate_sample.csv} al panel
\emph{Ingesta}. El sistema:

\begin{enumerate}
    \item Detecta el formato (FortiAnalyzer CSV) y carga un
    \code{FieldConfig} preconfigurado.
    \item Muestra preview con los roles de campo asignados
    (\code{TIME}, \code{ACTOR}, \code{TARGET}, \code{ACTION}).
    \item Pide confirmar o ajustar el mapping.
\end{enumerate}

Confirmar con los defaults. La ingesta tarda ~15-30 segundos para
el dataset demo. La barra de progreso muestra eventos procesados.

\section{3. Explorar el catálogo}

En el panel derecho, pestaña \emph{Catálogo}. Deberías ver:

\begin{itemize}
    \item \code{IP} — miles de direcciones observadas.
    \item \code{User} — cuentas con actividad.
    \item \code{Service} — apps identificadas por FortiOS.
\end{itemize}

Con conteos. Si algo se ve raro (conteos en 0, tipos inesperados),
revisá el field mapping — probablemente un rol quedó mal.

\section{4. Correr el primer hunt}

Abrir \emph{Hunt} (panel central). En la caja de hipótesis:

\begin{lstlisting}[basicstyle=\ttfamily\small]
IP -[Connect]-> IP {k=3}
\end{lstlisting}

Click \emph{Run}. El sistema devuelve paths ordenados por score.
Click en un path para ver su desglose:

\begin{itemize}
    \item \emph{Entity rarity}: 0.82 (las IPs son raras en conjunto).
    \item \emph{Edge rarity}: 0.91 (esta secuencia es inusual).
    \item \emph{Temporal novelty}: 0.64 (las IPs aparecieron tarde
    en la ventana).
\end{itemize}

\section{5. Modo explorador}

Hacer doble click sobre una IP interesante $\to$ abre el modo
\emph{explorador}. Ves los vecinos del nodo, las relaciones, la
timeline. Podés:

\begin{itemize}
    \item \code{Expand} para ver vecinos.
    \item \code{Expand grouped} si hay muchos, agrupados por tipo
    de relación.
    \item Ver \emph{eventos} del nodo en orden temporal.
\end{itemize}

\section{6. Etiquetar y anotar}

Sobre un nodo que te parece relevante, click derecho $\to$
\emph{Tag} $\to$ escribir ``sospechoso-c2''. Agregar una nota:
\emph{Note} $\to$ ``IP externa con patrón de beacon, ver
hunt \#2''.

\section{7. Exportar}

\emph{Archivo \textrightarrow{} Exportar \textrightarrow{} IOCs}.
Elegí formato STIX. Guardá el archivo. Ya tenés evidencia lista
para entregar.

\section{Y ahora?}

\begin{itemize}
    \item Entender los formatos: cap.~\ref{sec:man-ingesta}.
    \item Escribir hipótesis más elaboradas: cap.~\ref{sec:man-dsl}.
    \item Interpretar scoring: cap.~\ref{sec:man-caza}.
    \item Integrar con tu SIEM: cap.~\ref{sec:man-integraciones}.
\end{itemize}
